\documentclass[11pt,a4paper]{article}

% --- Packages ---
\usepackage[T1]{fontenc}
\usepackage[utf8]{inputenc}
\usepackage{lmodern}
\usepackage[margin=2.5cm]{geometry}
\usepackage{microtype}
\usepackage{setspace}
\usepackage{parskip}
\usepackage{booktabs}
\usepackage{xspace}
\usepackage{hyperref}
\hypersetup{colorlinks=true, linkcolor=blue, citecolor=blue, urlcolor=blue}

% --- Macros ---
\newcommand{\masld}{\textsc{MASLD}\xspace}
\newcommand{\metavir}{\textsc{METAVIR}\xspace}
\newcommand{\fibfour}{FIB-4\xspace}
\newcommand{\nfs}{\textsc{NFS}\xspace}
\newcommand{\apri}{\textsc{APRI}\xspace}
\newcommand{\te}{transient elastography\xspace}
\newcommand{\ipw}{\textsc{IPW}\xspace}
\newcommand{\strobe}{\textsc{STROBE}\xspace}
\newcommand{\pvalue}{$p$\xspace}
\newcommand{\Rsoftware}{\texttt{R~4.3.1}\xspace}
\newcommand{\survpkg}{\texttt{survival~3.5-7}\xspace}
\newcommand{\micepkg}{\texttt{mice~3.16.0}\xspace}

\begin{document}

\section*{Methods}

% [STROBE-4] Study design; [STROBE-5] Setting
This prospective observational cohort study was conducted across three tertiary hepatology
centres in mainland China (Beijing, Shanghai, and Guangzhou) and enrolled participants
between January 2018 and December 2022. The study was approved by the institutional review
boards of all three participating centres, and written informed consent was obtained from
every participant before enrolment. Reporting follows the \strobe statement for observational
cohort studies.

% [STROBE-6] Participants; [STROBE-7] Variables (exposure/outcome)
Eligible participants were adults aged 18 years or older with a diagnosis of \masld confirmed
by liver biopsy at baseline. Individuals were excluded if they had concurrent hepatitis~B or
hepatitis~C virus infection, reported alcohol consumption exceeding 30~g/day (men) or
20~g/day (women), had a prior or current diagnosis of any malignancy, or were pregnant at
the time of enrolment. Of 2{,}412 participants enrolled at baseline, 425 (17.6\%) were lost
to follow-up before the protocol-specified second biopsy, yielding a final analytic sample of
1{,}987 participants. The primary exposure comprised four non-invasive fibrosis scores
assessed at baseline: \fibfour, \nfs, \apri, and liver stiffness measured by \te
(FibroScan; Echosens, Paris, France) following the manufacturer's standardised protocol.
The primary outcome was fibrosis progression, defined as an increase of at least one stage
on the \metavir scale between the baseline and follow-up biopsies. Covariates specified
\textit{a priori} included age, sex, body-mass index, diabetes status, alcohol intake, and
baseline fibrosis stage.

% [STROBE-8] Data sources; [STROBE-9] Bias
Laboratory variables were obtained using harmonised protocols across all three centres to
minimise inter-site measurement variability. Liver biopsies were scored by local pathologists
using the \metavir system; to assess inter-rater reliability, a central pathologist
independently re-read a random 20\% sample (Cohen's $\kappa = 0.78$, indicating substantial
agreement). Three potential sources of bias were identified and addressed prospectively.
First, healthy-volunteer bias arising from selective refusal of the follow-up biopsy was
examined in a sensitivity analysis using inverse-probability weighting (\ipw), with weights
derived from logistic regression on baseline clinical characteristics. Second, inter-centre
misclassification of fibrosis stage was mitigated by the central re-read described above.
Third, residual confounding from unmeasured dietary factors was quantified using the
E-value approach.

% [STROBE-10] Study size
The target sample size was determined by a pre-specified power calculation: 1{,}800
participants were required to achieve 80\% power to detect a hazard ratio of 1.5 for
fibrosis progression (two-tailed $\alpha = 0.05$), assuming a 30\% three-year progression
rate. An additional 612 participants were recruited to allow for the anticipated 17--20\%
loss to follow-up, yielding the enrolled baseline cohort of 2{,}412.

% [STROBE-11] Statistical methods; [STROBE-12] Quantitative variables
The primary analysis used Cox proportional-hazards regression with time-varying covariates
to estimate the association between each baseline non-invasive score and the hazard of
fibrosis progression. All models were adjusted for age, sex, body-mass index, diabetes
status, alcohol intake, and baseline fibrosis stage. Missing covariate data were handled by
multiple imputation using chained equations ($m = 10$ imputed datasets; \micepkg). Estimates
were pooled across imputed datasets using Rubin's rules. Sensitivity analyses comprised
complete-case analysis and E-value calculation for unmeasured confounding. Three secondary
endpoints were pre-specified; the family-wise type~I error rate for these was controlled by
Bonferroni correction. All tests were two-tailed with a significance threshold of
$\alpha = 0.05$. Analyses were performed in \Rsoftware using the \survpkg and \micepkg
packages. Median follow-up was 3.4 years (interquartile range 2.6--4.1 years).

\end{document}
